\documentclass{article}

\usepackage{tabularx}
\usepackage{booktabs}

\title{Problem Statement and Goals\\\progname}

\author{\authname}

\date{}

\input{../Comments.text}
\input{../Common.text}

\begin{document}

\maketitle

\begin{table}[hp]
\caption{Revision History} \label{TblRevisionHistory}
\begin{tabularx}{\textwidth}{llX}
\toprule
\textbf{Date} & \textbf{Developer(s)} & \textbf{Change}\\
\midrule
Date1 & Name(s) & Description of changes\\
Date2 & Name(s) & Description of changes\\
... & ... & ...\\
\bottomrule
\end{tabularx}
\end{table}

\section{Problem Statement}

\wss{You should check your problem statement with the
\href{https://smiths.github.io/capTemplate/Checklists/ProbState-Checklist.pdf}
{problem statement checklist}}. 

\wss{You can change the section headings, as long as you include the required
information.}

\subsection{Problem}

\wss{What problem are you trying to solve?  Why is it important? Do not talk
about how you will solve the problem, only what the problem is.  AI is rarely
part of the problem.  The problem is ``what'' you want to do, not ``how'' you
want to do it.  The problem statement should be written in a way that is
understandable to a non-technical audience.}

\subsection{Inputs and Outputs}

\wss{Characterize the problem in terms of ``high level'' inputs and outputs.  
Use abstraction so that you can avoid details.}

\subsection{Stakeholders}

\subsection{Environment}

\wss{Hardware and Software Environment for the users.  The developer environment
is summarized as part of the developer plan.}

\section{Goals}

\section{Stretch Goals}

\wss{Goals you hope to get to, but not necessary.  They are also helpful if the
scope of the original project needs to be expanded.}

\section{Custom Documents (CDs)}

\wss{For CAS 741: State whether the project is a research project. This
designation, with the approval (or request) of the instructor, can be modified
over the course of the term.}

\wss{For SE Capstone: List your custom documents (CDs).  Potential CDs include
usability testing, code walkthroughs, user documentation, formal proof,
GenderMag personas, Design Thinking, etc.  (The full list is on the course
outline and in Lecture 02.) Normally the number of CDs will be two (one per
term).  Approval of the CDs will be part of the discussion with the instructor
for approving the project. The CDs, with the approval (or request) of the
instructor, can be modified over the course of the term.}

\textbf{Fall --- Literature Review Report}

We will review contemporary research on drawing practice and its relation to 
visual-motor development. We will also look the learning and support needs of 
children, including those with autism or ADHD. The literature review report 
will compare findings to identify limitations and explain how the evidence 
informs the application's activities, guidance, and progress tracking.

\textbf{Winter --- Usability Report}

We will evaluate the usability of our drawing activities and parent-facing 
progress summaries. The report will document evaluation methods and criteria 
as well as observed difficulties and recommended improvements. Such methods 
will include team walkthroughs with adult reviewers, not depending on child 
participation.

\newpage{}

\section*{Appendix --- Reflection}

\wss{Not required for CAS 741}

\input{../Reflection.text}

\subsection*{Shamil Canbolat}

\begin{enumerate}
    \item What went well while writing this deliverable? 

    The problem we chose to address already had research behind it with clearly 
    defined stakeholders, namely primary school teachers and children/students. 
    Our initial notes on the topic will translate easily into the literature 
    review custom document. Our supervisor and TA both have background knowledge 
    in education and thus directed us to clear and novel goals. Lastly, our team 
    is excited about the project, so we successfully assigned tasks, built a 
    workflow, and set up regular meetings with little conflict.

    \item What pain points did you experience during this deliverable, and how
    did you resolve them?

    Figuring out where to start with Gitlab/Avenue (checklists and rubrics) and 
    Github (issues, version control, and templates) was confusing at first, but 
    after coordinating online and in person and asking our TA during our meeting, 
    our team was on the same page. Also, attendance and participation was not the 
    best, but we outlined strict and detailed rules in the team charter which our 
    members enforced for productivity.

    \item How did you and your team adjust the scope of your goals to ensure
    they are suitable for a Capstone project (not overly ambitious but also of
    appropriate complexity for a senior design project)?

    We decided to limit our scope to just drawing practice for visual-motor 
    improvement, however, we will build our project with established design
    principles so that we can extend or pivot easily in the future. For example,
    the underlying structure of the drawing lessons can be duplicated for new 
    features like writing practice. Previous documentation and academic research 
    can be reused to implement AI features like drawing evaluation such that it 
    aligns with our original goals and standards. The final submission will be 
    fully useful with potential for development.

\end{enumerate}  

\end{document}